\documentclass{subfiles}
\begin{document}
    \begin{definition*}
        Let \(S\) be a set. We define a sequence of inclusion of the form
            \[
                S_{1} \subseteq S_{2} \subseteq \cdots \subseteq S_{k} \subseteq \cdots
            \]
            a (ascending) chain. Additionally, if there \(\exists n \in \N\) such 
            that \(S_{n} = S_{n + 1} = \cdots\), we say that the chain is stationary.
    \end{definition*}

    \begin{definition*}[Noetherian ring]
        Let \(R\) be a ring. We say that \(R\) is Noetherian if and only if 
            each ascending chain of ideals of \(R\) is stationary,
    \end{definition*}

    \begin{proposition*}
        Let \(R\) be a ring. Then, the following are equivalent.
        \begin{enumerate}
            \item \label{enm:1}  \(R\) is Noetherian.
            \item \label{enm:2} Each ideal \(I\) of \(R\) is finitely generated.
            \item \label{enm:3} However took a family \(\F*\) of ideals of \(R\); 
                \(\F*\) admits maximum. 
        \end{enumerate}
    \end{proposition*}

    \begin{proof*}
        \begin{description}
            \item [\(\ref{enm:2} \implies \ref{enm:1}\):] Let \(I_{1} \subseteq I_{2} \subseteq 
                \cdots \subseteq I_{k} \subseteq \cdots\) be a chain of ideals.
                Let  
                \[
                    I = \bigcup_{j = 1}^{\infty} I_{j}
                \]
                be an ideal finitely generated of \(R\).
                That is, \(I = (x_{1}, x_{2}, \ldots, x_{n})\).

                By the way \(I\) is defined, this means that 
                \[
                    \exists k_{1}, k_{2}, \ldots, k_{n} \text{ s.t. } x_{i} \in I_{k_{i}},
                \]
                hence, if \(m = \max{k_{i}}\) then \(x_{i} \in I_{m} = I_{m + 1} = \cdots\);
                that is, \(R\) is Noetherian.

            \item [\(\ref{enm:1} \implies \ref{enm:3}\):] Consider \(\F*\) be a non-empty family of ideals.
                Let \(I_{1} \in \F*\). It follows that 
                \begin{itemize}
                    \item \(I_{1}\) is maximal, from which the thesis.
                    \item \(I_{1}\) is not maximal. By hypothesis,
                        there exists \(I_{2} \supseteq I_{1}\).
                \end{itemize}
                Iterating the procedure we either find a maximal ideal,
                from which the thesis; or we end up with a non-stationary chain,
                which is absurd.

            \item [\(\ref{enm:3} \implies \ref{enm:3}\):] Let \(I \subseteq R\) be an ideal, 
                \(I \ne (0)\). Consider the following family of ideals
                \[
                    \F* = \set{J \subseteq R}[J \text{ is finitely generated}, J \subseteq I].
                \]
                Trivially, \(\F* \ne \varnothing\). Let \(I_{MAX}\) be the maximal element in \(\F*\).
                    Then \(I_{MAX} = (x_{1}, \ldots, x_{m})\) with \(x_{i} \in R\).
                    Assume, by the sake of contradiction that \(I_{MAX} \ne I\),
                    this implies that \(\exists y \in I\) such that \(y \notin I_{MAX}\).
                    But then we should have \((x_{1}, \ldots, x_{m}, y) \in \F*\),
                    meaning \(I_{MAX}\) isn't the maximal element, which is absurd.\qed
        \end{description}
    \end{proof*}

    \begin{proposition*}
        Let \(f : R \to Q\) be a ring homomorphism. Then, if \(f\) is onto and \(R\)
            is a Noetherian ring, \(Q\) is a Noetherian ring.
    \end{proposition*}

    \begin{proof*}
        Let \(J \subseteq Q\) be an ideal. By the properties of homomorphism
            \[
                I = f^{-1}(J) = \set*{a \in R}[f(a) \in J]
            \]
            is an ideal of \(R\).

        By hypothesis we have that \(R\) is Noetherian,
            then we can assume that \(I = (a_{1}, a_{2}, \ldots, a_{m})\).
            We prove that, \(J = (b_{1}, b_{2}, \ldots, b_{m})\), where \(b_{i} = f(a_{i})\).

        Consider any \(y \in J \implies \exists x \in I \text{ s.t. } f(x) = y\).
            But 
            \[
                x = \sum_{i = 0}^{n} a_{i} x_{i}, x_{i} \in R.
            \]
            Hence, we have that 
            \[
                \begin{aligned}
                    y & = f \left\lparen \sum_{i = 1}^{n} a_{i}x_{i} \right\rparen \\ 
                        & = \sum_{i = 1}^{n} f(a_{i}x_{i}) \\ 
                        & = \sum_{i = 1}^{n} \underbrace{f(a_{i})}_{= b_{i}}\underbrace{f(x_{i})}_{\in J}
                \end{aligned}
            \]
            therefore \(y \in (b_{1}, b_{2}, \ldots, b_{n})\) and \(J \subseteq (b_{1}, b_{2}, \ldots, b_{n})\).
            Proving that \((b_{1}, b_{2}, \ldots, b_{n}) \subseteq J\) is trivial.
    \end{proof*}

    \begin{theorem*}[Hilbert's base]
        Let \(R\) be a ring. If \(R\) is Noetherian, then \(R[x]\) is Noetherian.
    \end{theorem*}

    \begin{corollary*}
        Let \(\F\) be a field. Then, \(\F[x_{1}, x_{2}, \ldots, x_{n}]\) is Noetherian.
    \end{corollary*}

    \begin{proof*}
        Let us proceed by induction on the number of variables.
            By the previous theorem we know that for \(n = 1, \F[x_{1}]\) is Noetherian.
            As inductive step, assume the statement to be true for \(n - 1\);
            that is, \(\F[x_{1}, x_{2}, \ldots, x_{n - 1}]\) is Noetherian. 
            Then, we have that \(\F[x_{1}, x_{2}, \ldots, x_{n - 1}, x_{n}] = \F[x_{1}, x_{2}, \ldots, x_{n - 1}][x_{n}]\),
            which is true by hypothesis.\qed
    \end{proof*}

    \begin{corollary*}
        Given \(I \subseteq \F[x_{1}, x_{2}, \ldots, x_{n}]\), the ring 
        \[
            B = \frac{\F[x_{1}, x_{2}, \ldots, x_{n}]}{I}
        \]
        is Noetherian.
    \end{corollary*}

    \begin{proof*}
        Trivial. By assuming the epimorphism 
        \[
            \F[x_{1}, x_{2}, \ldots, x_{n}] \to \frac{\F[x_{1}, x_{2}, \ldots, x_{n}]}{I},
        \]
        the statement is true by the previously discussed proposition.\qed
    \end{proof*}

    \begin{proposition*}
        Let \(R\) be a Noetherian ring, and let \(J \subseteq R\) be an ideal.
            Then, \(\exists M \subseteq R\) maximal such that \(J \subseteq M\).
    \end{proposition*}

    \begin{proof*}
        Simply consider the family 
            \[
                \F* = \set*{I \subseteq R}[J \subseteq I].
            \]
            Since \(R\) is Noetherian, then \(\F*\) must admit a maximal element.\qed
    \end{proof*}
\end{document}
